% FACTORIZATION_ANCHORS.tex — VARIANT of Chapter 2.
%
% Identical to FACTORIZATION.tex except that assumption A1 is relaxed: converged objects are fed
% back to the room localiser as SE(2) landmarks, closing the loop the baseline chapter cuts.
% Diff against FACTORIZATION.tex to see exactly what changes.
%
% Build:  pdflatex FACTORIZATION.tex   (twice, for cross-references)
%
\documentclass[11pt,a4paper]{report}
\usepackage[T1]{fontenc}
\usepackage{lmodern}          % scalable Latin Modern — microtype needs a scalable face
\usepackage[margin=2.6cm]{geometry}
\usepackage{amsmath,amssymb,amsthm}
\usepackage{empheq}
\usepackage{pifont}
\usepackage{array}
\usepackage{booktabs}
\usepackage{enumitem}
\usepackage[hidelinks]{hyperref}
\usepackage{microtype}

\newcommand{\Ind}{\perp\!\!\!\perp}          % conditional independence
\newcommand{\given}{\,|\,}
\newcommand{\R}{\mathbb{R}}
\theoremstyle{definition}   % roman body: assumptions carry prose, not one-line claims
\newtheorem{assumption}{Assumption}
\renewcommand{\theassumption}{A\arabic{assumption}}   % A1 … A7, flat, not per-section

\begin{document}

\setcounter{chapter}{1}
\chapter{The generative model of the concept-agents\\[4pt]
         \large\normalfont\itshape variant: objects as localisation landmarks}

\begin{flushleft}
\textit{No\'e Zapata}
\end{flushleft}

\vspace{4pt}

\section{Introduction}
\label{sec:intro}



The \textsf{active\_inference} scene understanding runs as a set of concurrent processes: one
localiser, and one \emph{concept agent} per object class, each maintaining its own instances in a
shared graph. That split is usually described in engineering terms --- separation of
concerns, independent failure, parallelism.

It also admits a probabilistic reading, and that reading is the subject of this chapter. Taken
together, the processes compute an approximation to one posterior, and the way they are divided
\emph{is} a factorization of it. Each process boundary corresponds to a conditional-independence
assumption; each quantity that crosses a boundary is what survived a marginalisation; and several
behaviours that present as bugs turn out to be factors that were dropped deliberately and later
needed. Read this way, the architecture carries its failure modes on its face: they are readable
off the factors.

The chapter therefore writes that model down. It states the joint distribution over objects, room,
trajectory, association and sensor data (\S\ref{sec:joint}), conditions it on the data, says what
we are actually solving for (\S\ref{sec:problem}), and then states, one at a time, the seven
approximations that turn it into something the deployed processes compute
(\S\ref{sec:derivation}). Each is given with its justification, its cost, and the code that
implements it. Two of those costs are measured on the running system (\S\ref{sec:consequences}).

\medskip\noindent
\textbf{This is a variant chapter.} It differs from the baseline in one assumption, \ref{as:meanfield}.
There, the localiser is computed from odometry and LiDAR alone and the objects never inform it.
Here that channel is closed: a converged object is published back as an SE(2) landmark and enters
the localiser's own free energy. Everything else --- the joint, the conditioning, the other six
assumptions --- is unchanged, which is the point of keeping the two chapters parallel. The
consequences are concentrated in \S\ref{sec:derivation} and \S\ref{sec:fe}, and one of them is
serious enough to have kept the feature switched off (\S\ref{sec:incest}).

\medskip\noindent
One result is worth the effort on its own. The single exactly-solvable factor
(\S\ref{sec:exact}) produces a ceiling on how much one camera frame can ever contribute, which
accounts for a class of ``more data did not help'' outcomes. The lifecycle contract that every
agent must satisfy stands in a close relation to these seven assumptions, and is taken up
separately in the next chapter.

\section{Notation}
\label{sec:notation}

Everything below uses these, and nothing else. Every index set is finite.

\paragraph{Index sets.}
\begin{align*}
t &\in \{1,\dots,T\}    && \text{compute cycle (a mask frame with its capture stamp)}\\
c &\in \mathcal{C}       && \text{object \emph{class} — one concept agent per class:}\\
  &                      && \quad \mathcal{C}=\{\textsf{refrigerator},\textsf{table},\textsf{chair},\textsf{cabinet},\dots\}\\
i &\in \mathcal{I}_c     && \text{object \emph{instance} of class }c\text{ (a DSR node)}\\
j &\in \mathcal{J}_t     && \text{mask slice within frame }t\\
k &\in \mathcal{K}       && \text{\emph{burst} of measurements — \S\ref{sec:batch} only; unused in the deployed system}
\end{align*}

\paragraph{Unknowns.} What we are solving for, or must account for:
\begin{align*}
\theta_i &\in \R^{6}   && \text{geometry of instance }i:\ (c_x,c_y,H,w,h,\text{yaw})\\
e_i      &\in \R       && \text{existence of instance }i\text{, as log-odds}\\
M_t      &             && \text{\emph{matching} at cycle }t:\ \text{slice } j \mapsto \text{instance } i,\ \text{or } \varnothing\\
X        &=\{X_t\}     && \text{robot trajectory},\ X_t\in SE(2)\\
R        &             && \text{room geometry (the wall polygon)}
\end{align*}

\noindent
Five unknowns, not six: $\Delta=\{\delta_t\}$ appears throughout but is a \emph{coordinate} for
integrating $X$ out, not a degree of freedom of its own. It is defined in \S\ref{sec:deltaisx}
and used only in the per-frame marginalisation of \S\ref{sec:derivation}.

\noindent
$M$ ranges over assignments of slices to instances, one per cycle. It is \emph{not} an
association between accumulated groups of measurements and objects; no such variable exists here
(cf.\ \S\ref{sec:khronos}).

\paragraph{Given.}
\begin{align*}
Z &=\{Z_t\}       && \text{semantic masks from the camera, with 3-D support points}\\
L &=\{L_t\}       && \text{LiDAR sweeps}\\
\Phi &=\{\Phi_t\} && \text{odometry, } \Phi_t\sim (X_{t-1})^{-1}X_t
\end{align*}
We write $Z_t^{(i)}$ for the slice that $M_t$ assigns to instance $i$ (empty if none).

\paragraph{Covariances.} Three distinct objects, kept apart:
\begin{align*}
\Sigma_X^t  &&& \text{localization covariance of }X_t\text{, published by \textsf{room\_concept}}\\
\Sigma_\delta^t &&& \text{common-mode covariance of }\delta_t\ \text{(note: subscript }\delta,\ \text{not }c)\\
\Sigma_{\theta_i} &&& \text{posterior covariance of instance }i\text{'s geometry}
\end{align*}
with
\begin{equation}
\Sigma_\delta^t \;=\; \underbrace{\Sigma_{\mathrm{cm}}}_{\text{constant prior}}
\;+\; \underbrace{J\,\Sigma_X^t\,J^{\!\top}}_{\text{chain, capture-stamp pinned}} .
\label{eq:sigmadelta}
\end{equation}
$\Sigma_{\mathrm{cm}}$ is \texttt{common\_mode\_pos/size/yaw\_std}${}^2$; the second term is
\texttt{frame.chain\_cov\_*}. This is the \emph{only} place the foreign trajectory estimate
enters the object problem.

\section{The full joint, and where the conditional comes from}
\label{sec:joint}

\subsection{The generative model}

Before any conditioning, write the joint over \emph{everything} — unknowns and data alike — in
the order the world generates them. Each factor is a physical statement.

\begin{empheq}[box=\fbox]{align}
P\!\left(e,\theta,R,X,\Phi,L,M,Z\right)
\;=\;&
\underbrace{P(e)}_{\text{\ding{192} objects exist}}\;
\underbrace{P(\theta\given e)}_{\text{\ding{193} and have shape}}\;
\underbrace{P(R)}_{\text{\ding{194} the room has walls}}\;
\underbrace{P(X)}_{\text{\ding{195} robot moves}}
\nonumber\\[5pt]
\times\;&
\underbrace{P(\Phi\given X)}_{\text{\ding{196} wheels report it}}\;
\underbrace{P(L\given X,R,\theta)}_{\text{\ding{197} LiDAR hits walls \emph{and} objects}}
\nonumber\\[5pt]
\times\;&
\underbrace{P(M\given \theta,e,X)}_{\text{\ding{198} what is visible from where}}\;
\underbrace{P(Z\given \theta,M,X)}_{\text{\ding{199} surfaces become masks}}
\label{eq:fulljoint}
\end{empheq}

\noindent
Read as a story: objects exist \ding{192} and have geometry \ding{193}; the room has walls
\ding{194}; independently of both, the robot follows a trajectory \ding{195}, which the wheels
measure noisily \ding{196} and which the LiDAR measures far better by sweeping the walls
\ding{197}. From a given pose some objects are in the frustum and unoccluded, which is what a
matching can range over \ding{198}, and the visible surfaces project into mask points \ding{199}.

Three structural facts fall straight out of \eqref{eq:fulljoint}, and they matter more than the
algebra:

\begin{enumerate}[label=(\alph*),leftmargin=1.9em]
\item \textbf{Objects and pose are \emph{a priori} independent}: $P(e)P(\theta\given e)$ and
      $P(R)P(X)$ factorise exactly, with no assumption. Objects do not know where the robot is.
\item \textbf{They are coupled through \emph{two} channels, not one.} The obvious one is
      \ding{199}: masks are generated by object surfaces viewed from a pose. The less obvious one
      is \ding{197} --- \textbf{the LiDAR sweep contains the objects as well as the walls}. A
      refrigerator standing against a wall returns range in the very scan used to localise
      against that wall. In this variant \emph{both} bridges carry traffic: the mask bridge
      informs the objects, and the LiDAR bridge is supplemented by an explicit landmark channel
      that informs the localiser.
\item \textbf{$\Phi$ and $L$ are observations of $X$, not priors on it} --- factors \ding{196}
      and \ding{197}. They sit on the data side of the conditioning bar for that reason. Odometry
      ties consecutive poses together; the LiDAR, through $R$, is what stops the pose drifting
      without bound, and is therefore what makes $\Sigma_X$ small enough that
      \eqref{eq:sigmadelta} is dominated by its constant term.
\end{enumerate}

\subsection{$\Delta$ is a coordinate, not an unknown}
\label{sec:deltaisx}

Equation~\eqref{eq:fulljoint} has eight factors and no $\Delta$, and nothing below adds one. The
trajectory enters the problem exactly once, as $X$. What changes is the coordinate system in
which we eventually integrate it out.

That marginalisation is over a whole trajectory --- a large, strongly correlated object. But $X$
is never handed to us; what \textsf{room\_concept} publishes is an estimate $\hat X$ with
covariance $\Sigma_X$. It is therefore natural to integrate in the \emph{error} coordinate
\begin{equation}
\delta_t \;:=\; \hat X_t \ominus X_t \;\in\;\R^{6},
\qquad \Delta=\{\delta_t\}_{t=1}^{T},
\label{eq:deltadef}
\end{equation}
which, for the fixed known $\hat X$, is a bijection with unit Jacobian, so that for any $f$
\begin{equation}
\int f(X)\,\mathrm{d}X \;=\; \int f\!\left(\hat X\ominus\Delta\right)\,\mathrm{d}\Delta .
\label{eq:changeofvar}
\end{equation}

\begin{quote}
$\Delta$ carries no degree of freedom and appears in no factor of the model. It is a change of
integration variable, adopted for two reasons and no others: $\Sigma_X$ is expressed in it
(eq.~\ref{eq:sigmadelta}), and the \emph{per-frame} structure of the approximation
the per-frame approximation of \S\ref{sec:derivation} is visible in it while being invisible in
$X$. Treating $X$ and $\Delta$ as two
unknowns would double-count the same uncertainty.
\end{quote}

\noindent
Accordingly $\Delta$ does not appear again until \S\ref{sec:derivation}, where the integral is
actually performed. Everything between here and there is written in $X$.

\subsection{Conditioning}

The posterior is the joint over the evidence:
\begin{equation}
P\!\left(e,\theta,R,M,X\given Z,\Phi,L\right)
\;=\;
\frac{P\!\left(e,\theta,R,X,\Phi,L,M,Z\right)}{P(Z,\Phi,L)},
\qquad
P(Z,\Phi,L)=P(Z,L\given\Phi)\,P(\Phi).
\label{eq:conditioning}
\end{equation}
The denominator does not depend on any unknown, so it never affects the $\arg\max$ and is dropped
everywhere below. It is not inert, though: $-\log P(Z,\Phi,L)$ is the quantity the variational free
energy upper-bounds, so every assumption of \S\ref{sec:derivation} is paid for in exactly that
currency.

Substituting~\eqref{eq:fulljoint} into~\eqref{eq:conditioning} gives the posterior as an exact
equality --- the normaliser is written out, so nothing is hidden in a $\propto$:
\begin{align}
&P\!\left(e,\theta,R,M,X\given Z,\Phi,L\right)\;=\;\frac{1}{P(Z,\Phi,L)}\;\times
\nonumber\\[5pt]
&\qquad
\underbrace{P(e)\,P(\theta\given e)}_{\text{\ding{192}\ding{193}\; what objects are}}
\;\cdot\;
\underbrace{P(R)\,P(X)\,P(\Phi\given X)\,P(L\given X,R,\theta)}_{\text{\ding{194}--\ding{197}\; where the robot is}}
\nonumber\\[7pt]
&\qquad\times\;
\underbrace{P(M\given\theta,e,X)}_{\text{\ding{198}\; whose mask is whose}}
\;\cdot\;
\underbrace{P(Z\given\theta,M,X)}_{\text{\ding{199}\; how masks arise}}
\label{eq:posterior-raw}
\end{align}

\noindent
Everything so far is exact. \emph{No} approximation has been made: \eqref{eq:posterior-raw} is
just \eqref{eq:fulljoint} divided by its own marginal. The rest of this document consists of
replacing three of its four numerator blocks with something computable. Read as a roadmap:

\begin{description}[style=nextline,leftmargin=1.4em,itemsep=9pt,parsep=2pt]

\item[\normalfont \ding{192}\ding{193}\; $P(e)\,P(\theta\given e)$ \;---\; \textbf{what objects are}]
That a refrigerator exists at all, and that if it does, it is about $0.60\pm0.08$ deep and
$1.90\pm0.30$ tall. This is the one block we do \emph{not} approximate: these are the configured
priors, and the measurements in \S\ref{sec:consequences} argue they do more work than is usually
credited.

\item[\normalfont \ding{194}--\ding{197}\; $P(R)P(X)P(\Phi\given X)P(L\given X,R,\theta)$ \;---\; \textbf{where the robot is}]
Where the robot went and what the room looks like, constrained by the wheels, by the LiDAR
sweeping the walls, and --- in this variant --- by the converged objects themselves, fed back as
interior landmarks. Computed by \textsf{room\_concept}, but no longer independently of $\theta$:
the $\theta$ in the last factor is now honoured rather than dropped.

\item[\normalfont \ding{198}\; $P(M\given\theta,e,X)$ \;---\; \textbf{whose mask is whose}]
Properly a distribution over assignments: from this pose, with these objects, several matchings
are possible and some are likelier than others. Replaced by its single most likely value,
computed greedily.

\item[\normalfont \ding{199}\; $P(Z\given\theta,M,X)$ \;---\; \textbf{how masks arise}]
The second of the two bridges between objects and pose, and the one dealt with carefully.
Everything difficult lives here, and it is broken up along four different axes: across object
classes, across instances within a class, across time, and finally the trajectory is integrated
out of each frame separately.

\end{description}

\section{The problem}
\label{sec:problem}

What the agents have to publish is the geometry and existence of the objects: $\theta$ and $e$,
and nothing else. The room model, the robot trajectory and the mask-to-instance matching are all
things we are obliged to reason about --- the data cannot be explained without them --- but none
of them is a result. They are nuisance parameters.

The way to get from a distribution over everything to an answer about $\theta$ and $e$ is to
\emph{marginalise}: integrate the nuisance parameters out of the posterior~\eqref{eq:posterior-raw},
leaving a distribution over the two quantities we actually want, and take its maximum.

\begin{equation}
\boxed{\;
\theta^\star,\,e^\star
\;=\;
\underbrace{\arg\max_{\theta,\,e}}_{\text{maximised}}\;
\sum_{M}\;
\underbrace{\int}_{\text{marginalised}}
P\!\left(\theta,e,M,X,R \given Z,\Phi,L\right)\,
\mathrm{d}X\,\mathrm{d}R \; }
\label{eq:problem}
\end{equation}

\noindent
Written out, the marginalisation step alone is
\begin{equation}
P\!\left(\theta,e\given Z,\Phi,L\right)\;=\;\sum_M\iint
P\!\left(\theta,e,M,X,R\given Z,\Phi,L\right)\mathrm{d}X\,\mathrm{d}R ,
\label{eq:marginalpost}
\end{equation}
and \eqref{eq:problem} is its maximum. Note the structure here, since it is easy to read as a
posterior of a posterior, which it is not:
conditioning happened once, in \eqref{eq:conditioning}, and summing over $M$ and integrating over
$(X,R)$ merely removes variables from a distribution that is \emph{already} conditioned on
$(Z,\Phi,L)$. Both sides of \eqref{eq:marginalpost} are conditioned on the same thing.

What \eqref{eq:problem} maximises is therefore the \emph{marginal} posterior, not the joint one:
it is a marginal-MAP estimate. Maximising the joint instead --- reporting the $(\theta,e,M,X,R)$ that
jointly maximise~\eqref{eq:posterior-raw} --- is a different and generally non-equivalent
estimator, and it is the one a monolithic SLAM system computes.

Three treatments, three reasons:
\begin{center}
\begin{tabular}{@{}llp{8.1cm}@{}}
\toprule
Unknown & Treatment & Why \\
\midrule
$\theta,e$ & maximised & they are the answer; they are what we publish to DSR \\[2pt]
$X$, $R$   & marginalised & nuisance: needed to explain $Z$ and $L$, never reported by us. The
                            $X$ integral is carried out in the coordinate $\Delta$
                            (\S\ref{sec:deltaisx}) \\[2pt]
$M$        & \emph{collapsed} & neither maximised nor summed --- see \S\ref{sec:derivation} \\
\bottomrule
\end{tabular}
\end{center}

\noindent
The third row is the one to watch. Maximising and marginalising are both principled; collapsing
$M$ to its mode is neither, and it is the only place where a variable is treated as known when it
demonstrably is not. Its cost surfaces in \S\ref{sec:derivation}.

A system that reports its own trajectory --- a SLAM system --- would put $X$ in the first row
instead of the second. That single reassignment is what separates this factorization from a
monolithic one, and most of the derivation in \S\ref{sec:derivation} follows from it.

\section{Perception as free-energy minimisation}
\label{sec:fe}

Equation~\eqref{eq:posterior-raw} is exact and unusable: the marginal
\eqref{eq:marginalpost} has no closed form, and the joint has no representation an agent could
hold. What each agent holds instead is a Gaussian, and what it does each cycle is take a step
downhill on a bound.

\subsection{The bound}

For any distribution $q$ over the unknowns, the \emph{variational free energy}
\begin{equation}
F[q]\;=\;\mathbb{E}_{q}\!\left[-\ln P(\text{unknowns},\text{data})\right]-\mathbb{H}[q]
\;=\;\underbrace{\mathrm{KL}\!\left[\,q\,\|\,P(\text{unknowns}\given\text{data})\,\right]}_{\ge 0}
\;-\;\ln P(\text{data})
\label{eq:F}
\end{equation}
upper-bounds the surprise $-\ln P(Z,\Phi,L)$, with equality only when $q$ is the exact posterior.
Minimising $F$ over $q$ therefore does two things at once: it tightens the bound on the evidence,
and it drags $q$ towards the posterior we cannot compute. Everything the agents do to $\theta$ and
$e$ is an instance of this.

The choice of \emph{which} family $q$ is drawn from is not incidental --- it is the content of
\S\ref{sec:derivation}. Read as a statement about $q$ rather than about $P$, the seven assumptions
say:
\begin{equation}
q\!\left(\theta,e,M,X,R\right)\;=\;
\underbrace{q(X,R)}_{\text{foreign, but variable}}\;
\underbrace{\mathbb{1}\!\left[M=\hat M\right]}_{\text{degenerate}}\;
\prod_{c\in\mathcal{C}}\prod_{i\in\mathcal{I}_c}
\underbrace{q(\theta_i)}_{\mathcal{N}(\mu_i,\Sigma_{\theta_i})}\;
\underbrace{q(e_i)}_{\text{log-odds}} .
\label{eq:qfamily}
\end{equation}
Each factor of \eqref{eq:qfamily} that is narrower than the truth adds a non-negative amount to the
KL term of \eqref{eq:F}, and that amount is the cost quoted with each assumption.

\subsection{What each block of the posterior contributes}

Substituting \eqref{eq:qfamily} into \eqref{eq:F} and using the four blocks of
\eqref{eq:posterior-raw}:

\begin{center}
\begin{tabular}{@{}l l p{6.9cm}@{}}
\toprule
Block of \eqref{eq:posterior-raw} & Contributes to $F$ & \\
\midrule
\ding{192}\ding{193} objects &
$\mathrm{KL}\!\left[q(\theta_i)\|P(\theta_i\given e_i)\right]$ &
the \emph{complexity} term: how far this instance's posterior has been dragged from the
class prior. Minimised by us. \\[4pt]
\ding{194}--\ding{197} localisation &
$\mathbb{E}_{q}\!\left[-\ln P(L\given X,R,\theta)\right]$ &
no longer inert. The object's pose enters the localiser's accuracy term, so moving
$\theta$ moves this block too --- which is what the landmark channel transmits. \\[4pt]
\ding{198} matching &
zero &
a delta contributes no entropy and no expectation. The matching is free, in the
literal sense that $F$ cannot see it. \\[4pt]
\ding{199} masks &
$\mathbb{E}_{q}\!\left[-\ln P(\bar Z_i\given\theta_i,\hat X)\right]$ &
the \emph{accuracy} term. All the data enter here. Minimised by us. \\
\bottomrule
\end{tabular}
\end{center}

\noindent
Only the matching block is inert now. Each side still minimises a \emph{partial} free energy
holding the other fixed --- for an object agent,
\begin{equation}
F_i\;=\;
\underbrace{\mathbb{E}_{q}\!\left[-\ln P\!\left(\bar Z_i\given\theta_i,\hat X\right)\right]}_{\text{accuracy}}
\;+\;
\underbrace{\mathrm{KL}\!\left[\,q(\theta_i)\,\|\,P(\theta_i\given e_i)\,\right]}_{\text{complexity}} ,
\label{eq:freeenergy}
\end{equation}
summed over instances within a class --- and, by \ref{as:class}, never summed across classes. The
localiser minimises its own partial free energy over $(X,R)$, in which \eqref{eq:anchormsg} appears
as one factor per anchor. Alternating the two is coordinate descent on their sum: each step lowers
the joint free energy, and the pair converges to a fixed point of the factorised problem.

The accuracy term deserves one remark, because it is where the robustness lives. It is not a
least-squares residual over the mask points but the negative log of a \emph{mixture}: each point is
explained either by the object's surface, by a uniform clutter component, or by a known wall. A
point that fits nothing costs a bounded amount and stops pulling, so an outlier degrades the fit by
a finite quantity instead of dragging it. That bound is a modelled probability, not a rejection
threshold.

\subsection{How the minimisation is carried out}

$q(\theta_i)$ is Gaussian, so minimising \eqref{eq:freeenergy} reduces to finding its mode and
curvature. The engine takes a few Gauss--Newton steps on $F_i$, then adopts the Hessian at the
optimum as $\Sigma_{\theta_i}$ --- which is what makes the result a \emph{Laplace} posterior rather
than a point estimate. Between observations the transition prior of \ref{as:markov} is applied on
its own, so an unobserved instance loses precision rather than holding a stale certainty.

\subsection{The same quantity scores action}

Perception minimises $F$ with respect to $q$. The other half of the active-inference formulation
minimises \emph{expected} free energy with respect to what the robot does next, and the agents
implement its epistemic term directly: the next-best-view planner scores each face of the fitted
box by the D-optimal expected entropy reduction on $\Sigma_{\theta_i}$, and proposes the standoff
pose that observes the winner. The covariance produced by the Hessian in the previous paragraph is
the same object the planner differentiates to choose where to look. Belief and behaviour are scored
in one currency.

\subsection{What free energy is not used for}

One decision is conspicuously outside this scheme. Whether an instance should exist at all is a
comparison between models of different order, and it has a free-energy form --- but no such
comparison is computed. Birth is decided by an accumulated weight in different units, and the
mismatch is taken up in the next chapter.

\section{The derivation}
\label{sec:derivation}

Start from the exact posterior~\eqref{eq:posterior-raw}. Each assumption below replaces one of its
numerator blocks; none of them touches the denominator. There are seven, and this is the first
place they are named:

\begin{center}
\begin{tabular}{@{}>{\bfseries}l l p{7.6cm}@{}}
\toprule
\normalfont Assumption & Acts on & What it assumes \\
\midrule
\ref{as:meanfield} & \ding{194}--\ding{197} & localisation and objects are updated alternately, each holding the other fixed \\
\ref{as:frames}    & \ding{199} & the trajectory marginal factorises frame by frame \\
\ref{as:hard}      & \ding{198} & one matching, not a distribution over matchings \\
\ref{as:class}     & \ding{199} & object classes are independent --- one agent each \\
\ref{as:inst}      & \ding{199} & instances within a class are independent \\
\ref{as:geomexist} & \ding{199} & geometry and existence factorise \\
\ref{as:markov}    & \ding{199} & the state is Markov in time \\
\bottomrule
\end{tabular}
\label{tbl:assumptions}
\end{center}

\noindent
Five of the seven act on \ding{199} alone. That concentration is expected: \ding{199} is where
objects and pose meet, and it is the only block large enough to need breaking up. Block \ding{192}\ding{193} ---
the object priors --- appears nowhere in the table, because nothing is done to it.

\begin{assumption}[Structured exchange — objects as landmarks]\label{as:meanfield}
\begin{equation}
q\!\left(\theta,e,X,R\right)\;=\;q(X,R)\;\prod_{c}\prod_{i}q(\theta_i)\,q(e_i),
\end{equation}
with both factors \emph{variable}. Neither is computed from the other in one pass; they are
updated alternately, each holding the other fixed. The posterior itself is unchanged --- the
objects were always inside $P(L\given X,R,\theta)$ --- and what is dropped is only the
\emph{correlation} between the two blocks, not the dependence.

\smallskip
This is coordinate descent on the joint free energy, and it is a strictly weaker assumption than a
one-way cut: each block still assumes the other is a point-mass-with-covariance rather than a
correlated partner, but information now flows in both directions.
\end{assumption}

\noindent\textit{Justification.} Walls constrain pose well in a closed room and badly in two
common situations: a long corridor, where position along the axis is weakly observable, and an open
area where the nearest wall is far enough that its angular resolution is poor. Furniture sits in the
interior. A converged, well-localised object is exactly the landmark those geometries lack, and the
information is already being computed --- \S\ref{sec:derivation} produces $\Sigma_{\theta_i}$ for
every instance whether or not anyone consumes it.

\smallskip\noindent\textit{Cost.} One is structural and one is dangerous.
\begin{enumerate}[label=(\roman*),leftmargin=1.9em,topsep=3pt,itemsep=3pt]
\item The mean-field form still discards $\mathrm{Cov}(\theta,X)$, so both blocks remain
      overconfident. Alternating updates converge to a fixed point of the factorised problem, not
      to the true joint posterior.
\item Errors become \emph{correlated across agents}. A mis-localised fridge now perturbs $X$, which
      perturbs every chair, table and cabinet in the room. \ref{as:inst} assumes instances own
      disjoint measurement sets and are therefore independent; with the loop closed they share a
      common pose, and that independence is no longer even approximately true. The cost of one bad
      object stops being local.
\end{enumerate}

\smallskip\noindent\textit{Computed as.} Two messages, one per direction.

\emph{Localiser $\to$ objects}, as before: $(\hat X,\Sigma_X)$, consumed through
\texttt{InnerGaussianAPI} at the mask's capture stamp and folded into $\Sigma_\delta^t$ by
\eqref{eq:sigmadelta}.

\emph{Objects $\to$ localiser}: each agent publishes, per cycle, its instance's pose in the
\emph{robot} frame together with an information matrix,
\begin{equation}
z_o\;=\;\big(x,y,\mathrm{yaw}\big)_{\text{robot}},
\qquad
\Lambda\;=\;\big(\Sigma_o\oplus R_o\big)^{-1},
\label{eq:anchormsg}
\end{equation}
where $\Sigma_o$ is the object's own map covariance and $R_o$ this frame's measurement covariance.
The localiser adds one factor per anchor to its window. Three guards are already in place: an
anchor is only emitted once its map covariance is below a validation bound, its weight decays with
staleness, and only the newest few window slots receive object factors. A symmetric object
contributes position only --- the yaw row and column of $\Lambda$ are zeroed rather than the yaw
residual being trusted.

The machinery for this exists in full: the object side publishes today
(\texttt{PublishObjectObs}), and the room side carries the matching factor. What is off is the
room-side switch, for the reason in \S\ref{sec:incest}.

\subsection{Why this is not simply switched on}
\label{sec:incest}

The object's belief was formed \emph{using} $\hat X$. Sending $z_o$ back to the localiser, which
produces $\hat X'$, which updates the object, circulates the same evidence: the LiDAR sweep that
localised the robot also carried the returns from the fridge, and the mask that fitted the fridge
was de-projected with the pose that sweep produced. Iterating a mean-field exchange over shared
evidence makes both blocks progressively more confident without any new observation --- the
estimate becomes consistent with itself rather than with the room.

The formal requirement is that $\Lambda$ in \eqref{eq:anchormsg} carry only the \emph{innovation}:
the information the object contributes beyond what it received. The published form,
$(\Sigma_o\oplus R_o)^{-1}$, is an attempt at that --- inflating by this frame's measurement
covariance before inverting --- but it is not a marginalisation of the prior that came from
$\hat X$, and under repeated exchange the residual overlap accumulates.

Doing this properly is a known problem with known answers, and they differ in what they choose to
represent. The next section takes them in turn.

\subsection{Handling the correlation properly}
\label{sec:correlation}

Three treatments, and a fourth that mixes them.

\paragraph{(a) Represent it: one joint Gaussian.} Stop factorising. Let
$\xi=(X,R,\theta_1,\dots,\theta_n)$ and keep $q(\xi)=\mathcal{N}(\mu,\Sigma)$ with dense
cross-blocks. The free energy \eqref{eq:F} is unchanged, but its KL term no longer decomposes:
there is no per-instance $F_i$ to hand to an agent, only one objective and one Gauss--Newton step
over the whole state. In practice one works in information form, because $\Lambda=\Sigma^{-1}$ is
sparse exactly where factors couple variables --- a mask couples $(X_t,\theta_i)$, a sweep couples
$(X_t,R)$, odometry couples $(X_t,X_{t-1})$ --- and the system is solved by sparse Cholesky. This
is graph SLAM with object landmarks.

Double counting cannot occur: each measurement appears once, as one factor, and is never re-sent.
Loop closure deforms the objects automatically. What it costs is the architecture. \ref{as:class}
and \ref{as:inst} both dissolve --- once the information matrix couples classes through shared
poses, no single process can own a class --- and bounding the window requires marginalising old
poses, which produces \emph{fill-in}: every landmark seen from a marginalised pose becomes directly
correlated with every other. The usual repairs (sparsification, generic linear constraints, a
Chow--Liu approximation of the dense block) are themselves approximations, so one approximation is
traded for another.

\paragraph{(b) Keep the factorisation, make the messages honest.} Retain
$q(X,R)\prod_i q(\theta_i)$ and repair the \emph{messages} instead. Before instance $i$ sends its
anchor to the localiser, divide out the message it sent previously, so what it transmits is only
the information it has added since. In natural parameters this is a subtraction, so it is cheap.
This is the cavity construction of expectation propagation, and it removes the pathology of
\S\ref{sec:incest} directly: evidence can no longer circulate, because no factor is ever counted
twice.

The process separation survives intact --- each agent still owns its block, and the protocol still
carries only $(\hat X,\Sigma_X)$ one way and $(z_o,\Lambda)$ the other. What it does \emph{not}
recover is the correlation itself: both marginals remain overconfident, they simply stop becoming
progressively more so. The new cost is bookkeeping: each agent must store its outstanding message
per consumer, which is bounded and small.

\paragraph{(c) Assume nothing: fuse conservatively.} Covariance intersection combines two estimates
of unknown correlation as
\begin{equation}
\Lambda_{\text{ci}}\;=\;\omega\,\Lambda_a+(1-\omega)\,\Lambda_b,
\qquad \omega\in[0,1],
\end{equation}
with $\omega$ chosen to minimise the trace or determinant of the result. The fused estimate is
guaranteed never optimistic, whatever the true correlation, and it requires no provenance, no
cross-covariance and no change of protocol. The price is tightness: information is discarded to
stay safe. Since the motivation for anchors is precisely to \emph{tighten} the pose where walls
constrain it poorly, conservative fusion may return much of what the anchors were introduced to
buy --- a real tension, not a detail.

\paragraph{(d) A structured middle.} The correlation that matters is between an object and the
handful of poses it was actually observed from, over a short interval. Keeping a dense block over
$(X_{t-k:t},\theta_i)$ while factorising across everything else is a \emph{structured} mean field:
exact where the dependence is strong, factorised where it is weak. The bounded-window idea appears
here applied to the correlation rather than to the measurements.

\medskip\noindent
Of the four, (a) is the correct answer and the wrong project: it replaces the concept-agent
architecture with a monolithic estimator, and the chapter it belongs to is a different one. (c) is
safe and may be self-defeating for the reason just given. \textbf{(b) is the one that fits} --- it
preserves every process boundary, adds one stored message per consumer and nothing else, and
removes exactly the failure that keeps the room-side switch off today. It is the prerequisite this
variant should acquire before being enabled.

\begin{assumption}[The marginalisation over $X$ factorises per frame]\label{as:frames}
Carrying out $\int(\cdot)\,q(X)\,\mathrm{d}X$ in the coordinate $\Delta$ of~\eqref{eq:changeofvar},
\begin{equation}
\int P\!\left(Z\given\theta,M,X\right)q(X)\,\mathrm{d}X
\;\approx\;
\prod_t \int P\!\left(Z_t\given\theta,M,\hat X_t\ominus\delta_t\right)
\mathcal{N}\!\left(\delta_t;\,0,\,\Sigma_\delta^t\right)\mathrm{d}\delta_t
\label{eq:perframe}
\end{equation}
with $\Sigma_\delta^t$ as in~\eqref{eq:sigmadelta}. This is the \emph{only} place $\Delta$ is
used, and the first place it was needed.
\end{assumption}
\noindent\textit{Justification.} Each frame is de-projected with its own capture-stamp transform,
so its offset is drawn afresh. \textit{Cost.} \textbf{False across short windows.} Localisation
error does not resample every $100$\,ms; over a fraction of a second the true $\delta_t$ are
strongly correlated, and over a short enough window they are near-constant. Sequential per-frame
updates never expose this, because each frame is consumed alone and its offset integrated out
immediately. Any procedure that \emph{pools} several frames' raw points into one estimate does
expose it, and breaks --- measured in \S\ref{sec:consequences}(b).


\smallskip\noindent\textit{Computed as.} By never instantiating $\delta_t$. The integral in \eqref{eq:perframe} is Gaussian, so Woodbury
turns it into a cap on how much a single frame may contribute: \texttt{recursive\_laplace.h} sets
$\Sigma_c^{-1}\leftarrow$ \texttt{m.common\_mode\_inv\_diag(frame)} and folds it into the
information sum, which is \eqref{eq:sat}. The entire cost of the trajectory's uncertainty is thus
one $6$-vector per frame, and no pose variable is ever touched by the object agent.

\begin{assumption}[Hard matching]\label{as:hard}
\begin{equation}
P(M\given X,Z,\Phi)\;\approx\;\mathbb{1}\!\left[M=\hat M\right],
\qquad
\hat M=\arg\max_M P(M\given \hat X,Z)
\end{equation}
computed greedily, by lowest squared Mahalanobis distance over gated pairs, in the shared
\texttt{InstanceTracker} with gate $\chi^2_2=9$.
\end{assumption}
\noindent\textit{Justification.} Tractability only. \textit{Cost.} The sum over $M$
in~\eqref{eq:problem} collapses to one term, so a wrong match is unrecoverable within the cycle;
the merge operator and the existence channel are the repair, and they act after the damage.


\smallskip\noindent\textit{Computed as.} By \texttt{rc::InstanceTracker::update()}. For every (detection, instance) pair it forms the
innovation covariance $S=P+R^2I$ from the instance's own posterior $\Sigma_{\theta_i}$, gates on
squared Mahalanobis distance $\le\chi^2_2=9$, and assigns greedily, lowest cost first, one-to-one.
$\hat M$ then does nothing but route slice indices onto instances. No cost of the matching enters
$F_i$ --- which is exactly why a mis-assignment is invisible to the free energy that would
otherwise have caught it.

\begin{assumption}[Class partition]\label{as:class}
\begin{equation}
P(\theta,e\given \hat M,X,Z,\Phi)\;=\;\prod_{c\in\mathcal{C}}
P\!\left(\theta_{\mathcal{I}_c},e_{\mathcal{I}_c}\given \hat M,X,Z,\Phi\right)
\end{equation}
\end{assumption}
\noindent\textit{Justification.} Architectural: one agent per class, and the label is fixed by
which binary you are, not inferred. \textit{Cost.} The sharpest hole in the whole scheme. A
cabinet mis-labelled \textsf{refrigerator} cannot be resolved by comparing the two classes'
likelihoods, because \textsf{refrigerator\_concept} and \textsf{cabinet\_concept} never see each
other's. \texttt{apply\_fridge\_filter} (singleton inhibition $+$ shape-plausibility decay) is a
\emph{within-class} substitute for what should be across-class model comparison.


\smallskip\noindent\textit{Computed as.} By running one process per class. Every concept agent subscribes to the same \texttt{"masks"}
node in the graph and keeps only the slices whose label matches its own, so each minimises its own
$\sum_i F_i$ over a disjoint set of measurements. No term anywhere compares the free energy a
cabinet hypothesis would achieve against the one a refrigerator hypothesis achieves: both numbers
exist, in two different processes, and are never brought together.

\begin{assumption}[Instance independence given the matching]\label{as:inst}
\begin{equation}
\prod_{c}P(\theta_{\mathcal{I}_c},e_{\mathcal{I}_c}\given\cdot)
\;=\;\prod_{c}\prod_{i\in\mathcal{I}_c}P(\theta_i,e_i\given\cdot)
\end{equation}
\end{assumption}
\noindent\textit{Justification.} $\hat M$ is one-to-one, so instances own disjoint measurement
sets. \textit{Cost.} Exact given \ref{as:hard} \emph{except} for physical exclusion — two
fridges cannot occupy one volume — which is a genuine coupling reintroduced procedurally by
\texttt{merge\_overlapping\_instances}.


\smallskip\noindent\textit{Computed as.} This is where the whole scheme actually runs, so it is
worth following in full. Each cycle, \texttt{process\_refrigerator\_node()} is called once per
live instance, and the body is:

\begin{quote}
\ttfamily\small
for (idx : inst.assigned\_mask\_idxs)\\
\hspace*{2em}obs = observe\_slice(inst, idx);\\
\hspace*{2em}F = run\_inference(inst, obs);
\end{quote}

\noindent
Four properties of that loop are modelling decisions, though they read as implementation
detail.

\emph{One update per slice, not one per instance.} If a fridge is seen by two sensors in the same
cycle, the loop runs twice. Sequential Bayesian updates compose into the joint likelihood, so this
is exact --- but the reason it is done this way rather than by concatenating both clouds into one
frame is \ref{as:frames}: each slice carries its \emph{own} $\Sigma_\delta^t$, because two sensors
do not share a registration error. Concatenation would force one common-mode term on both and
silently over-count their agreement.

\emph{The instance must earn its update.} The loop is entered only when the mask packet is fresh
in \emph{two} senses: a new \texttt{frame\_id} \emph{and} a new capture timestamp. The publish
counter alone is not enough --- it advances on every republish even when the camera is frozen, and
re-integrating one capture as if it were independent evidence makes the belief ratchet: geometry
rotates and reshapes with both robot and scene stationary. The second gate is the difference
between $T$ genuine observations and one observation counted $T$ times, which is exactly the
distinction $\prod_t$ in \eqref{eq:result} presumes.

\emph{Classify, don't destroy.} \texttt{observe\_slice()} splits the assigned cloud by signed
distance into points that lie on the current surface and points that do not. Both are kept: the
first drive the fit, the second are the residual that later evidence --- a second face, a
neighbouring object, a mis-association --- may claim. Nothing is discarded on the strength of the
current estimate.

\emph{No cross terms, by construction.} $F$ for the class is the plain sum $\sum_i F_i$. Two
instances that overlap in space produce no penalty in the free energy at all, because their
likelihoods are computed over disjoint measurement sets and never compared. Physical exclusion ---
two refrigerators cannot occupy one volume --- is therefore restored \emph{outside} the model, by
an operator that deletes one of any two instances whose oriented footprints overlap. That operator
is the visible scar of this assumption: it exists because \eqref{eq:result} has nowhere to put the
constraint.

When no slice was assigned, the loop body never runs and only the predict step of
\ref{as:markov} fires, so an unobserved instance ages rather than freezes.

\begin{assumption}[Geometry/existence mean-field]\label{as:geomexist}
\begin{equation}
P(\theta_i,e_i\given\cdot)\;\approx\;q(\theta_i)\,q(e_i)
\end{equation}
\end{assumption}
\noindent\textit{Justification.} None, strictly — this one is convenience. \textit{Cost.} Shape
and existence \emph{are} coupled: an implausible shape is evidence of non-existence. The coupling
survives as a one-way channel (plausibility $\rightarrow$ existence log-odds), not as a joint.


\smallskip\noindent\textit{Computed as.} By keeping two separate objects on the instance: \texttt{ai2\_belief}, the Gaussian over
geometry, and an existence log-odds accumulated from LiDAR occupancy carve and mask-silhouette
evidence. They meet in one direction only --- the fitted shape's plausibility becomes a bounded
log-evidence ratio that is added to the existence channel, so a badly-shaped object decays. The
reverse, existence tightening the geometry, has no path.

\begin{assumption}[Markov time, conditionally independent frames]\label{as:markov}
\begin{equation}
q(\theta_i)\;\propto\;\prod_t
\Big[\textstyle\int P\!\left(Z_t^{(i)}\given\theta_i^t,\hat X_t\ominus\delta_t\right)
\mathcal{N}\!\left(\delta_t;0,\Sigma_\delta^t\right)\mathrm{d}\delta_t\Big]\;
P\!\left(\theta_i^t\given\theta_i^{t-1}\right)
\end{equation}
\end{assumption}
\noindent\textit{Justification.} Rigid static furniture; small process noise. \textit{Cost.}
Negligible here. Note what this factor commits us to: estimation is \emph{recursive and
immediate}. Each measurement is folded in as it arrives and then discarded, so there is no point
at which a decision can be deferred until more evidence is in. That is a real design choice with
real consequences, measured in \S\ref{sec:consequences}.


\smallskip\noindent\textit{Computed as.} By the \texttt{predict()} half of the recursive-Laplace filter: the mean is held and
$\Sigma_{\theta_i}\!\leftarrow\!\Sigma_{\theta_i}+Q$ with a small per-frame $Q$. The
\texttt{update()} half then takes a few Gauss--Newton steps on \eqref{eq:freeenergy} and adopts
the Hessian at the optimum as the new $\Sigma_{\theta_i}$ --- which is what makes the posterior a
\emph{Laplace} one. On a stale cycle only \texttt{predict()} runs, on the agent's own clock, so a
dead sensor shows up as growing covariance instead of unchanged confidence.

\subsection*{The assembled result}

Substituting \ref{as:meanfield}--\ref{as:markov} into~\eqref{eq:posterior-raw} and marginalising $X$
by~\eqref{eq:perframe} leaves a posterior over the objects and the matching alone:

\begin{empheq}[box=\fbox]{align}
P\!\left(\theta,e,M\given Z,\Phi\right)
\;\propto\;
&\underbrace{\mathbb{1}\!\left[M=\hat M\right]}_{\text{\ref{as:hard}}}\;
\prod_{c\in\mathcal{C}}\;\prod_{i\in\mathcal{I}_c}
\underbrace{q(e_i)}_{\text{\ref{as:geomexist}}}
\nonumber\\[2pt]
\times\;\prod_t\;
&\underbrace{\int P\!\left(Z_t^{(i)}\given\theta_i^t,\hat X_t\ominus\delta_t\right)
\mathcal{N}\!\left(\delta_t;0,\Sigma_\delta^t\right)\mathrm{d}\delta_t}_{\text{\ref{as:frames}}}
\;\underbrace{P\!\left(\theta_i^t\given\theta_i^{t-1}\right)}_{\text{\ref{as:markov}}}
\label{eq:result}
\end{empheq}

\noindent
Note where each assumption ended up. Three are \emph{visible} as factors
(\ref{as:hard}, \ref{as:geomexist}, \ref{as:markov}) and one as the integral
(\ref{as:frames}). The other three act \emph{structurally} and leave no factor of their own:
\ref{as:class} and \ref{as:inst} are the two products $\prod_c\prod_i$, and
\ref{as:meanfield} --- the trajectory --- survives only inside $\Sigma_\delta^t$, having been
integrated away. That last point is worth dwelling on: the foreign estimate does not appear in
the object posterior as a factor at all. It appears as a \emph{covariance}, and nowhere else.
Every symbol in~\eqref{eq:result} was introduced in \S\ref{sec:notation}.

\section{The one factor that is exact}
\label{sec:exact}

Only the $\delta_t$ integral is done without approximation (given a linearisation). It is worth
seeing why, because it produces the ceiling that governs everything downstream.

For frame $t$, stack the $N$ point residuals of one mask. Writing $G\in\R^{N\times 6}$ for the
Jacobian of those residuals with respect to the shared offset, and $\varepsilon\sim
\mathcal{N}(0,\sigma^2 I)$ for the independent per-point noise,
\begin{equation}
r \;=\; r_0 + G\,\delta_t + \varepsilon,
\qquad \delta_t\sim\mathcal{N}(0,\Sigma_\delta^t).
\end{equation}
Marginalising $\delta_t$ gives $r\sim\mathcal{N}\!\left(r_0,\ \sigma^2 I+G\Sigma_\delta^t G^{\!\top}\right)$,
so the frame's Fisher information about $\theta_i$ is
\begin{equation}
\mathcal{I}_t \;=\; G^{\!\top}\!\left(\sigma^2 I+G\Sigma_\delta^t G^{\!\top}\right)^{-1}\!G
\;\overset{\text{Woodbury}}{=}\;
\left(\sigma^2\left(G^{\!\top}G\right)^{-1}+\Sigma_\delta^t\right)^{-1}
\;\xrightarrow[\ N\to\infty\ ]{}\;
\left(\Sigma_\delta^t\right)^{-1}.
\label{eq:sat}
\end{equation}
using $A-A(B+A)^{-1}A=(A^{-1}+B^{-1})^{-1}$ with $A=G^{\!\top}G$, $B=\sigma^2(\Sigma_\delta^t)^{-1}$.

Two readings of~\eqref{eq:sat}, both load-bearing:
\begin{enumerate}[leftmargin=1.6em]
\item \textbf{A denser mask stops helping.} Information saturates at $(\Sigma_\delta^t)^{-1}$.
      That ceiling --- exactly $1/(\sigma^2_{\mathrm{cm}}+\Sigma_{\text{chain}})$ per degree of
      freedom --- is what \texttt{common\_\allowbreak mode\_\allowbreak inv\_\allowbreak diag()}
      returns.
\item \textbf{$\delta_t$ is never computed, because it cannot be.} The residual depends on
      $\theta_i$ and $\delta_t$ only through their sum, so the joint information for
      $(\theta_i,\delta_t)$ has rank $6$, not $12$. Only the ceiling is computable. What
      \emph{is} estimated is $\hat X$ — that is what localisation is; $\delta_t$ is its residual
      error, and if it were knowable it would already have been removed.
\end{enumerate}

\clearpage
\appendix
\renewcommand{\thesection}{\Alph{section}}
\setcounter{section}{0}

\begin{center}
{\Large\bfseries Appendices}\\[4pt]
\rule{0.32\textwidth}{0.4pt}
\end{center}

\medskip\noindent
The body above is the model and its derivation. What follows is everything empirical or
comparative: two measurements that test the assumptions (\S\ref{sec:consequences}), the shape a
batch layer would have to take if the model ever outgrew them (\S\ref{sec:batch}), and the
correspondence with the formulation this one was derived against (\S\ref{sec:khronos}).

\section{Two consequences, both measured}
\label{sec:consequences}

Both follow from factors of~\eqref{eq:result}, and both were checked on the running system rather
than argued.

\paragraph{(a) A batch layer would buy nothing at this model complexity.}
A recurring proposal is to insert an intermediate layer between measurements and objects: hold a
burst of measurements while they accumulate, commit to an object only at the end. The ceiling
in~\eqref{eq:sat} says when that pays. It binds only when $\theta_i$ has more degrees of freedom
than one frame can determine. A non-parametric surface --- a mesh or a TSDF, thousands of DOF, no
shape prior --- is far past that line and \emph{must} accumulate. Our $\theta_i$ is six
parameters with tight priors ($0.60\pm0.08$ footprint, $1.90\pm0.30$ height), and it is not:
live, an instance born from a single $8\,000$-point mask starts at $w=0.669$, $h=0.600$ and
converges to $w=0.674$, $h=0.604$. The priors already do the work the accumulation was meant to
do. This is a property of the \emph{model}, not of the tuning.

\paragraph{(b) Pooling frames violates \ref{as:frames}, and it shows.}
Feeding a burst of $8$ frames' raw points to one geometric moment asks the data to separate eight
different $\delta_t$ from geometry --- impossible by the identifiability argument of \S\ref{sec:exact}(2), and with no per-burst
offset for them to land in. Live A/B on the same object:

\begin{center}
\begin{tabular}{@{}lccc@{}}
\toprule
& single frame & pooled burst & converged \\
\midrule
$w$ at birth        & $0.669$ & $\mathbf{0.915}$ & $0.673$ \\
$h$ at birth        & $0.600$ (prior) & $\mathbf{0.494}$ & $0.626$ \\
$\sigma_H$ at birth & $0.300$ & $0.190$ & --- \\
cycles to settle $w$ & $\sim 0$ & $\sim 255$ & --- \\
\bottomrule
\end{tabular}
\end{center}

\noindent
Note the third row: the pooled estimate was \emph{more confident} about a worse answer, which is
the mean-field signature of \ref{as:meanfield} compounding with \ref{as:frames}. The diagnostic
that separates the two hypotheses is that \emph{both} footprint extents inflated: a genuine
second face of the object raises the minor axis and leaves the major axis at the true width,
whereas pose-error blur stretches both. Reverted; see the $\star$ note in
\texttt{RefrigeratorFitter::run\_inference}.

\section{When a batch layer would be worth adding}
\label{sec:batch}

Should the model ever outgrow \S\ref{sec:consequences}(a), \eqref{eq:result} says exactly what the
extra layer must carry. Group the measurements into bursts $\bar Z_k$, $k\in\mathcal{K}$, and give
each burst its own pose $T_k$ --- a nuisance at \emph{burst} granularity rather than frame
granularity:
\begin{equation}
\prod_{k\in\mathcal{K}} \int P\!\left(\bar Z_k\given\theta_i,T_k\right)\,P\!\left(T_k\right)\,\mathrm{d}T_k .
\label{eq:burstnuisance}
\end{equation}
Three requirements, none currently met:
\begin{enumerate}[leftmargin=1.6em]
\item bank the burst in \emph{camera} frame and chain it internally by odometry. Only the
      relative pose inside a short window is trustworthy; absolute room-frame coordinates bake in
      $\hat X(t)$ and cannot be un-baked
      (\texttt{MaskIngestor::enable\_frame\_transform} already provides this path);
\item lift the common-mode nuisance from per-frame to per-burst, i.e.\ give $T_k$ a prior and
      integrate it out, exactly as \S\ref{sec:exact} does for the per-frame offset $\delta_t$;
\item supply something that makes $T_k$ identifiable --- a co-observed static anchor, or a
      revisit of the same object from a different viewpoint. Without it, an unobservable
      parameter has merely been added, and the identifiability argument of \S\ref{sec:exact}(2)
      applies to it verbatim.
\end{enumerate}
Worth doing only when $\theta_i$ gains degrees of freedom a single frame cannot determine: a
non-parametric surface, an articulated object, an open-set detection with no shape prior. For a
six-DOF box with priors, no.


\section{Relation to Khronos}
\label{sec:khronos}

For readers coming from Khronos \cite{khronos}, whose eq.~(16) is
\begin{equation}
P(O,X,Y,A\given Z,\Phi)=
\underbrace{\prod_i P(O_i\given \bar Y_i,X)}_{\text{reconciliation}}\;
\underbrace{P(X,A\given Y,\Phi)}_{\text{SLAM}}\;
\underbrace{\prod_k P(Y_k\given \bar Z_k,\bar\Phi_k)}_{\text{local estimation}} ,
\end{equation}
the correspondence with~\eqref{eq:result} is:

\begin{center}
\begin{tabular}{@{}p{4.2cm}p{2.2cm}p{7.4cm}@{}}
\toprule
Khronos & Here & Consequence \\
\midrule
$\prod_k P(Y_k\given\bar Z_k,\bar\Phi_k)$ \newline \emph{object fragments} &
absent &
Their $Y_k$ is the batch layer of \S\ref{sec:batch}, with $T_{RY_k}$ as the burst pose
of~\eqref{eq:burstnuisance}. Necessary for them because $\Omega_k$ is a mesh; not for us
(\S\ref{sec:consequences}(a)). \\[4pt]
$P(X,A\given Y,\Phi)$ &
$q(X)$ only &
They maximise the trajectory jointly with the map, so a loop closure deforms object poses
retroactively. \ref{as:meanfield} forecloses that. \\[4pt]
$A:\ Y\mapsto \bar Y_i$ &
\textbf{no counterpart} &
$A$ associates \emph{fragments to objects}; with no $Y$ there is nothing for it to range over.
Our $M_t$ is instead their \emph{in-window} $Z\!\to\!Y$ step, solved greedily by volumetric IoU
and not named in eq.~(16) --- it lives inside the local-estimation term. \\[4pt]
$\prod_i P(O_i\given\bar Y_i,X)$ &
$q(e_i)$ &
Existence runs live here, from occupancy-carve and silhouette evidence, rather than being
recovered afterwards by geometric verification. \\[4pt]
--- & $\prod_{c}$ &
The class partition \ref{as:class} has no analogue there; neither does its arbitration hole. \\
\bottomrule
\end{tabular}
\end{center}

\noindent
Their local-consistency premise (eq.~8) --- that the \emph{relative} pose error inside a window is
bounded --- is the same statement as the failure mode of \ref{as:frames}: it is what licenses
pooling in camera frame, and what forbids pooling in room frame.

\begin{thebibliography}{9}
\bibitem{khronos}
L.~Schmid, M.~Abate, Y.~Chang and L.~Carlone.
\emph{Khronos: A Unified Approach for Spatio-Temporal Metric-Semantic SLAM in Dynamic
Environments.} Proc. Robotics: Science and Systems, 2024. arXiv:2402.13817.
\end{thebibliography}

\end{document}
